\documentclass[11pt]{article}

\usepackage[margin=1in]{geometry}
\usepackage{amsmath,amssymb,amsthm}
\usepackage{enumitem}
\usepackage{graphicx}
\usepackage{hyperref}

% Common PDE notation
\newcommand{\R}{\mathbb{R}}
\newcommand{\pd}[2]{\frac{\partial #1}{\partial #2}}
\newcommand{\pdd}[2]{\frac{\partial^2 #1}{\partial #2^2}}
\newcommand{\lap}{\Delta}
\newcommand{\grad}{\nabla}

% One numbered problem per \begin{problem}...\end{problem}
\theoremstyle{definition}
\newtheorem{problem}{Problem}

\title{APM5557: Partial Differential Equations\\Homework Problems}
\author{Frank Palazzolo}
\date{}

\begin{document}

\maketitle

\begin{problem}
From page 40 in the notes.

for the PDE: 
$u_t - c^2 u_{xx} = 0$, 
with initial condition $u(x,0) = g(x)$, 
show that the solution is given by the formula:
\begin{equation} \label{eq:solution}
u(x,t) = \frac{1}{2c\sqrt{\pi t}} \int_{-\infty}^{\infty} e^{-\frac{(x-\xi)^2}{4c^2t}} g(\xi) d\xi
\end{equation}
\end{problem}

\begin{proof}[Solution]
First compute the time derivative of $u$. Bringing $\pd{}{t}$ inside the
integral gives
\begin{equation} \label{eq:ut-inside}
u_t = \int_{-\infty}^{\infty} \pd{}{t}\left[
  \frac{1}{2c\sqrt{\pi t}} \, e^{-\frac{(x-\xi)^2}{4c^2t}}
\right] g(\xi) \, d\xi.
\end{equation}
Applying the product rule to the terms in brackets,
\begin{equation} \label{eq:ut}
u_t = \int_{-\infty}^{\infty} \left[
  -\frac{1}{4c\sqrt{\pi}\, t^{3/2}} \, e^{-\frac{(x-\xi)^2}{4c^2t}}
  + \frac{1}{2c\sqrt{\pi t}} \cdot \frac{(x-\xi)^2}{4c^2t^2} \,
    e^{-\frac{(x-\xi)^2}{4c^2t}}
\right] g(\xi) \, d\xi.
\end{equation}
Next compute the spatial derivative of $u$ in the same way. Bringing
$\pd{}{x}$ inside the integral gives
\begin{equation} \label{eq:ux-inside}
u_x = \int_{-\infty}^{\infty} \pd{}{x}\left[
  \frac{1}{2c\sqrt{\pi t}} \, e^{-\frac{(x-\xi)^2}{4c^2t}}
\right] g(\xi) \, d\xi.
\end{equation}
Since the prefactor does not depend on $x$, only the exponential is
differentiated, so
\begin{equation} \label{eq:ux}
u_x = \int_{-\infty}^{\infty}
  \frac{1}{2c\sqrt{\pi t}} \cdot \left(-\frac{x-\xi}{2c^2t}\right)
  e^{-\frac{(x-\xi)^2}{4c^2t}} g(\xi) \, d\xi.
\end{equation}
Taking one more derivative of \eqref{eq:ux} with respect to $x$ gives
\begin{equation} \label{eq:uxx-inside}
u_{xx} = \int_{-\infty}^{\infty}
  \frac{1}{2c\sqrt{\pi t}} \, \pd{}{x}\left[
    -\frac{x-\xi}{2c^2t} \, e^{-\frac{(x-\xi)^2}{4c^2t}}
  \right] g(\xi) \, d\xi.
\end{equation}
Applying the product rule to the terms in brackets,
\begin{equation} \label{eq:uxx}
u_{xx} = \int_{-\infty}^{\infty}
  \frac{1}{2c\sqrt{\pi t}} \left[
    -\frac{1}{2c^2t} \, e^{-\frac{(x-\xi)^2}{4c^2t}}
    + \frac{(x-\xi)^2}{4c^4t^2} \, e^{-\frac{(x-\xi)^2}{4c^2t}}
  \right] g(\xi) \, d\xi.
\end{equation}
Multiplying \eqref{eq:uxx} by $c^2$ cancels some of the factors of $c$, so
\begin{equation} \label{eq:c2uxx}
c^2 u_{xx} = \int_{-\infty}^{\infty}
  \frac{1}{2c\sqrt{\pi t}} \left[
    -\frac{1}{2t} \, e^{-\frac{(x-\xi)^2}{4c^2t}}
    + \frac{(x-\xi)^2}{4c^2t^2} \, e^{-\frac{(x-\xi)^2}{4c^2t}}
  \right] g(\xi) \, d\xi.
\end{equation}
Substituting \eqref{eq:ut} and \eqref{eq:c2uxx} into the PDE gives
\begin{equation} \label{eq:subst}
\begin{split}
u_t - c^2 u_{xx}
  &= \int_{-\infty}^{\infty} \left[
    -\frac{1}{4c\sqrt{\pi}\, t^{3/2}} \, e^{-\frac{(x-\xi)^2}{4c^2t}}
    + \frac{1}{2c\sqrt{\pi t}} \cdot \frac{(x-\xi)^2}{4c^2t^2} \,
      e^{-\frac{(x-\xi)^2}{4c^2t}}
  \right] g(\xi) \, d\xi \\
  &\quad - \int_{-\infty}^{\infty}
  \frac{1}{2c\sqrt{\pi t}} \left[
    -\frac{1}{2t} \, e^{-\frac{(x-\xi)^2}{4c^2t}}
    + \frac{(x-\xi)^2}{4c^2t^2} \, e^{-\frac{(x-\xi)^2}{4c^2t}}
  \right] g(\xi) \, d\xi.
\end{split}
\end{equation}
Since
\begin{equation} \label{eq:prefactor}
\frac{1}{2c\sqrt{\pi t}} \cdot \frac{1}{2t}
  = \frac{1}{4c\sqrt{\pi}\, t^{3/2}},
\end{equation}
the two integrals in \eqref{eq:subst} have identical integrands, so
everything cancels and
\begin{equation} \label{eq:pde}
u_t - c^2 u_{xx} = 0.
\end{equation}
\end{proof}

\begin{problem}
From 3.3-1 in the book.

Let $\Omega$ be a smooth bounded domain in $\mathbb{R}^n$. 
For a $C^2$ solution of $u$ of the wave equation $u_{tt} = c^2 \lap{u}$, 
show that if the energy is defined as $E_\Omega(t) = \int_\Omega (u_t^2 + c^2 |\grad u|^2) dx$, 
and either $u(x,t) = 0$ or $\frac{\partial u}{\partial n} = 0$ on $\partial\Omega$, 
then $E_\Omega(t)$ is constant in time.
\end{problem}

\begin{proof}[Solution]
To show that $E_\Omega(t)$ is constant, take its time derivative and set it
equal to 0:
\begin{equation} \label{eq:dEdt}
\frac{dE_\Omega}{dt} = \frac{d}{dt} \, \frac{1}{2} \int_\Omega
  \left( u_t^2 + c^2 |\grad u|^2 \right) dx = 0.
\end{equation}
Splitting this into two integrals,
\begin{equation} \label{eq:dEdt-split}
\frac{d}{dt} \, \frac{1}{2} \int_\Omega u_t^2 \, dx
  + \frac{d}{dt} \, \frac{1}{2} \int_\Omega c^2 |\grad u|^2 \, dx = 0.
\end{equation}
Bringing $\frac{d}{dt}$ inside the first integral, $u_t^2$ becomes
$2 u_{tt} u_t$, which cancels the $\frac{1}{2}$ for that term:
\begin{equation} \label{eq:dEdt-first}
\int_\Omega u_{tt} u_t \, dx
  + \frac{d}{dt} \, \frac{1}{2} \int_\Omega c^2 |\grad u|^2 \, dx = 0.
\end{equation}
Bringing $\frac{d}{dt}$ inside the second integral and writing
$|\grad u|^2 = \sum_{i=1}^n u_{x_i}^2$, each $u_{x_i}^2$ becomes
$2 u_{x_i} u_{x_i t}$, which cancels the $\frac{1}{2}$ for that term. Since
$c$ is constant, we can also pull $c^2$ out of the integral:
\begin{equation} \label{eq:dEdt-second}
\int_\Omega u_{tt} u_t \, dx
  + c^2 \int_\Omega \sum_{i=1}^n u_{x_i} u_{x_i t} \, dx = 0.
\end{equation}
Next we show that the second integral in \eqref{eq:dEdt-second} is equivalent
to the integral of $-u_t \lap u$, by integrating by parts once in each $x_i$.
Since $u_{x_i t} = (u_t)_{x_i}$, each term of the sum is
$u_{x_i} (u_t)_{x_i}$. Integrating by parts in $x_i$ moves the derivative off
of $u_t$ and onto $u_{x_i}$. By the divergence theorem, with $n_i$ the $i$-th
component of the outward normal,
\begin{equation} \label{eq:ibp-component}
\int_\Omega u_{x_i} (u_t)_{x_i} \, dx
  = \int_{\partial\Omega} u_t \, u_{x_i} n_i \, dS
  - \int_\Omega u_t \, u_{x_i x_i} \, dx.
\end{equation}
Summing over $i$, and using $\sum_{i=1}^n u_{x_i} n_i = \grad u \cdot n
= \pd{u}{n}$ and $\sum_{i=1}^n u_{x_i x_i} = \lap u$,
\begin{equation} \label{eq:ibp-sum}
\int_\Omega \sum_{i=1}^n u_{x_i} u_{x_i t} \, dx
  = \int_{\partial\Omega} u_t \, \pd{u}{n} \, dS
  - \int_\Omega u_t \lap u \, dx.
\end{equation}
(This is Green's first identity with $f = u_t$ and $g = u$.)
The boundary integral vanishes under either boundary condition. If $u = 0$ on
$\partial\Omega$ for all $t$, then $u$ does not change in time there, so
$u_t = 0$ on $\partial\Omega$. If $\pd{u}{n} = 0$ on $\partial\Omega$, then
that factor is zero directly. Therefore
\begin{equation} \label{eq:ibp-result}
\int_\Omega \sum_{i=1}^n u_{x_i} u_{x_i t} \, dx
  = -\int_\Omega u_t \lap u \, dx.
\end{equation}
Substituting \eqref{eq:ibp-result} into \eqref{eq:dEdt-second},
\begin{equation} \label{eq:dEdt-lap}
\int_\Omega u_{tt} u_t \, dx - c^2 \int_\Omega u_t \lap u \, dx = 0.
\end{equation}
Combining the terms and factoring out $u_t$,
\begin{equation} \label{eq:dEdt-factored}
\frac{dE_\Omega}{dt}
  = \int_\Omega u_t \left( u_{tt} - c^2 \lap u \right) dx = 0.
\end{equation}
The second factor is the original PDE, which is 0. Since the change in energy
is 0, $E_\Omega(t)$ is constant in time.
\end{proof}

\begin{problem}
From 3.3-2a in the book.

For the non-homogeneous wave equation $u_{tt} = \lap u + f(x,t)$,
show uniqueness of the solution $u$ given initial conditions $u(x,0) = g(x)$
\end{problem}

\begin{proof}
    Assume we have two solutions $u_1$ and $u_2$ to the non-homogeneous wave equation with the same initial conditions. Define $u = u_1 - u_2$. Then $u$ satisfies the homogeneous wave equation:   

$u_{tt} = \lap u$

$u = 0$ on $\partial \Omega$

By the previous problem, we know that the energy $E(t) = \int_\Omega (u_t^2 + |\grad u|^2) dx$ is constant in time.

\end{proof}

\begin{problem}
From 3.3-2b in the book.

For the non-homogeneous wave equation $u_{tt} = \lap u + f(x,t)$,
show uniqueness of the solution $u$ given initial conditions $u_t(x,0) = g(x)$
\end{problem}

\begin{proof}
    Assume we have two solutions $u_1$ and $u_2$ to the non-homogeneous wave equation with the same initial conditions. Define $u = u_1 - u_2$. Then $u$ satisfies the homogeneous wave equation:   

$u_{tt} = \lap u$

$u_t = 0$ on $\partial \Omega$

By the previous problem, we know that the energy $E(t) = \int_\Omega (u_t^2 + |\grad u|^2) dx$ is constant in time.

\end{proof}

\begin{problem}
From 3.3-6 in the book.

\end{problem}

\begin{proof}

We will use Green's first identity, which states that for a $C^1$ scalar
field $f$ and a $C^2$ scalar field $g$ on $\overline{\Omega}$,
\begin{equation} \label{eq:green}
\int_\Omega \left( f \lap g + \grad f \cdot \grad g \right) dx
  = \int_{\partial\Omega} f \, \pd{g}{n} \, dS.
\end{equation}
Substituting $u$ for both $f$ and $g$ in \eqref{eq:green}, and using
$\grad u \cdot \grad u = |\grad u|^2$,
\begin{equation} \label{eq:green-u}
\int_\Omega \left( u \lap u + |\grad u|^2 \right) dx
  = \int_{\partial\Omega} u \, \pd{u}{n} \, dS.
\end{equation}
If either $u(x,t) = 0$ or $\pd{u}{n} = 0$ on $\partial\Omega$, then the
integrand on the right side of \eqref{eq:green-u} vanishes, so that term is
zero and
\begin{equation} \label{eq:green-zero}
\int_\Omega \left( u \lap u + |\grad u|^2 \right) dx = 0.
\end{equation}
Rearranging \eqref{eq:green-zero},
\begin{equation} \label{eq:grad-lap}
\int_\Omega |\grad u|^2 \, dx = \int_\Omega -u \lap u \, dx.
\end{equation}

\end{proof}

\end{document}
